\section*{Problems}

\subsection*{Problem statements}

% Remember
\begin{problema}
	\label{prob:b94325d}
	Let $S$ be the sample space and let $A, B \subset S$ be two subsets. Explicitly list the $2^2 = 4$ disjoint elements that make up the fundamental partition $\particion(A,B)$ of the space $S$.
\end{problema}

% Understand
\begin{problema}
	\label{prob:320d62d}
	Extend the construction of a partition to the case of three subsets $A, B, C \subset S$. How many elements does the partition $\particion(A,B,C)$ have, and what are these elementary disjoint subsets explicitly? Explain why the number of elements doubles when a third set is added.
\end{problema}

% Apply
\begin{problema}
	\label{prob:18eada9}
	In a group of 50 students, 35 play soccer, 25 play basketball, and 15 play both sports. Using the variables of the fundamental partition $\particion(F,B)$, how many students do not play either sport?
\end{problema}

% Analyze
\begin{problema}
	\label{prob:6d1bd44}
	Using the principle of inclusion-exclusion or partition theory, prove the cardinality identity for the union of three sets $A, B, C \subset S$:
	\begin{align}
		\card{A \cup B \cup C} = \card{A} + \card{B} + \card{C} - \card{A \cap B} - \card{A \cap C} - \card{B \cap C} + \card{A \cap B \cap C}.
	\end{align}
\end{problema}

% Evaluate
\begin{problema}
	\label{prob:15ead08}
	A team of analysts in an epidemiological clinical trial with a total of $\card{S} = 500$ patients evaluated the presence of three symptoms of a new disease ($S_1, S_2, S_3$). The preliminary report presents the following census statistics:
	\begin{itemize}
		\item $\card{S_1} = 300, \quad \card{S_2} = 250, \quad \card{S_3} = 200$
		\item $\card{S_1 \cap S_2} = 180, \quad \card{S_1 \cap S_3} = 150, \quad \card{S_2 \cap S_3} = 130$
		\item $\card{S_1 \cap S_2 \cap S_3} = 30$
	\end{itemize}
	Using the algebraic breakdown into the 8 elements of the partition $\particion(S_1,S_2,S_3)$, rigorously evaluate why the report contains an unavoidable logical error that makes it impossible for these data to correspond to a real census. Identify specifically which partition element yields a mathematically absurd cardinality.
\end{problema}

% Create
\begin{problema}
	\label{prob:95d5458}
	A product team for a music-streaming app classifies its 800 active users according to three consumption habits: listening to \emph{podcasts} ($P$), creating their own playlists ($L$), and using offline mode ($O$). Design an original example with hypothetical cardinalities for $\card{P}$, $\card{L}$, $\card{O}$ and their intersections (single, double, and triple) such that the data are logically consistent (unlike the clinical-trial problem). Compute, using the breakdown into the 8 elements of $\particion(P,L,O)$, how many users have none of the three habits.
\end{problema}

\subsection*{Detailed solutions}

% Remember
\begin{solproblema}[prob:b94325d]
	For two subsets $A, B \subset S$, each element of the space $S$ may belong or not belong (complement) to each of the two sets. This generates $2 \times 2 = 4$ mutually disjoint elementary intersections whose union is the whole space $S$:
	\begin{itemize}
		\item $E_0 = A' \cap B'$ (elements outside both sets)
		\item $E_1 = A' \cap B$ (elements belonging only to $B$)
		\item $E_2 = A \cap B'$ (elements belonging only to $A$)
		\item $E_3 = A \cap B$ (elements in the common intersection of both sets)
	\end{itemize}
	By construction, $E_i \cap E_j = \emptyset$ for every $i \neq j$ and $\bigcup_{k=0}^3 E_k = S$.
\end{solproblema}

% Understand
\begin{solproblema}[prob:320d62d]
	For three subsets $A, B, C \subset S$, including binary membership with respect to a third set $C$ or its complement $C'$ doubles the number of regions relative to the two-set case: $2^3 = 8$ disjoint elementary elements.

	These 8 subsets, ordered systematically in binary code (where belonging to the set turns a bit on and belonging to the complement turns it off), are:
	\begin{itemize}
		\item $E_0 = A' \cap B' \cap C'$ (none of the three sets)
		\item $E_1 = A' \cap B' \cap C$ (only in $C$)
		\item $E_2 = A' \cap B \cap C'$ (only in $B$)
		\item $E_3 = A' \cap B \cap C$ (simultaneously in $B$ and $C$, but not in $A$)
		\item $E_4 = A \cap B' \cap C'$ (only in $A$)
		\item $E_5 = A \cap B' \cap C$ (simultaneously in $A$ and $C$, but not in $B$)
		\item $E_6 = A \cap B \cap C'$ (simultaneously in $A$ and $B$, but not in $C$)
		\item $E_7 = A \cap B \cap C$ (in the simultaneous intersection of the three sets)
	\end{itemize}
	The number doubles because each new set adds one independent binary bit (belongs / does not belong) to the coding of each region: moving from 2 to 3 sets multiplies the number of possible combinations by 2, from $2^2=4$ to $2^3=8$.
\end{solproblema}

% Apply
\begin{solproblema}[prob:18eada9]
	Let $F$ be the set of students who play soccer and $B$ the set of students who play basketball. From the partition of the sample space into four disjoint regions, define the elementary cardinalities:
	\begin{align}
		x_3 &= \card{F \cap B} = 15 \\
		x_2 &= \card{F \backslash B} = \card{F} - \card{F \cap B} = 35 - 15 = 20 \\
		x_1 &= \card{B \backslash F} = \card{B} - \card{F \cap B} = 25 - 15 = 10
	\end{align}
	The total number of students who play at least one sport is the union $\card{F \cup B} = x_1 + x_2 + x_3 = 10 + 20 + 15 = 45$.

	By complementarity with respect to the 50 students in the group:
	\begin{align}
		x_0 &= \card{(F \cup B)'} = \card{S} - \card{F \cup B} = 50 - 45 = 5.
	\end{align}
	Exactly 5 students do not play either sport.
\end{solproblema}

% Analyze
\begin{solproblema}[prob:6d1bd44]
	First consider the addition theorem for two sets $A, B$, with the cardinalities of the partition $\particion(A,B)$: $x_1 = \card{A' \cap B}$, $x_2 = \card{A \cap B'}$, $x_3 = \card{A \cap B}$. Since the elementary sets of a partition are mutually disjoint, the cardinality of any union is the direct sum of the cardinalities of its component regions, which gives $\card{A} + \card{B} - \card{A \cap B} = \card{A \cup B}$.

	For three sets, express the union by grouping two terms and applying the previous theorem:
	\begin{align}
		\card{A \cup B \cup C} &= \card{(A \cup B) \cup C} = \card{A \cup B} + \card{C} - \card{(A \cup B) \cap C}.
	\end{align}
	By the distributive law of intersection over union, $(A \cup B) \cap C = (A \cap C) \cup (B \cap C)$. Applying the two-set formula again to this distributed intersection:
	\begin{align}
		\card{(A \cup B) \cap C} &= \card{A \cap C} + \card{B \cap C} - \card{A \cap B \cap C}.
	\end{align}
	Substituting this and $\card{A \cup B} = \card{A} + \card{B} - \card{A \cap B}$ into the first equation:
	\begin{align}
		\card{A \cup B \cup C} &= \card{A} + \card{B} + \card{C} - \card{A \cap B} - \card{A \cap C} - \card{B \cap C} + \card{A \cap B \cap C}.
	\end{align}
	This proves algebraically the inclusion-exclusion formula of order three.
\end{solproblema}

% Evaluate
\begin{solproblema}[prob:15ead08]
	To verify the logical and mathematical consistency of the report, compute the variables $x_k = \card{E_k}$ ($k=0,1,\dots,7$) for the partition generated by the three symptoms ($S_1, S_2, S_3$):
	\begin{align}
		x_7 &= \card{S_1 \cap S_2 \cap S_3} = 30 \\
		x_6 &= \card{S_1 \cap S_2 \cap S_3'} = \card{S_1 \cap S_2} - x_7 = 180 - 30 = 150 \\
		x_5 &= \card{S_1 \cap S_2' \cap S_3} = \card{S_1 \cap S_3} - x_7 = 150 - 30 = 120 \\
		x_3 &= \card{S_1' \cap S_2 \cap S_3} = \card{S_2 \cap S_3} - x_7 = 130 - 30 = 100
	\end{align}
	Now infer the number of patients with only the first symptom ($x_4 = \card{S_1 \cap S_2' \cap S_3'}$):
	\begin{align}
		\card{S_1} = x_7 + x_6 + x_5 + x_4 \implies 300 = 30 + 150 + 120 + x_4 \implies x_4 = 0 \text{ (no contradiction so far).}
	\end{align}
	Now infer the number of patients with only the second symptom ($x_2 = \card{S_1' \cap S_2 \cap S_3'}$):
	\begin{align}
		\card{S_2} = x_7 + x_6 + x_3 + x_2 \implies 250 = 30 + 150 + 100 + x_2 \implies x_2 = -30.
	\end{align}
	\textbf{Conclusion:} the calculation yields a negative cardinality ($x_2 = -30$) for the region of patients who present only symptom $S_2$. In the axioms of counting and set theory, the cardinality function $\card{\cdot} \ge 0$ does not admit negative values. This proves beyond doubt that the figures published in the preliminary report are mutually inconsistent and reflect an error in the census collection or tabulation of intersections.
\end{solproblema}

% Create
\begin{solproblema}[prob:95d5458]
	\textbf{Example:} let $\card{S}=800$, with $\card{P}=350$ (podcasts), $\card{L}=400$ (own playlists), $\card{O}=300$ (offline mode), $\card{P \cap L}=150$, $\card{P \cap O}=120$, $\card{L \cap O}=180$, and $\card{P \cap L \cap O}=60$.

	Compute the 8 regions of $\particion(P,L,O)$:
	\begin{align}
		x_7 &= \card{P \cap L \cap O} = 60 \\
		x_6 &= \card{P \cap L \cap O'} = \card{P \cap L} - x_7 = 150-60 = 90 \\
		x_5 &= \card{P \cap L' \cap O} = \card{P \cap O} - x_7 = 120-60 = 60 \\
		x_3 &= \card{P' \cap L \cap O} = \card{L \cap O} - x_7 = 180-60 = 120 \\
		x_4 &= \card{P \cap L' \cap O'} = \card{P} - (x_7+x_6+x_5) = 350-(60+90+60) = 140 \\
		x_2 &= \card{P' \cap L \cap O'} = \card{L} - (x_7+x_6+x_3) = 400-(60+90+120) = 130 \\
		x_1 &= \card{P' \cap L' \cap O} = \card{O} - (x_7+x_5+x_3) = 300-(60+60+120) = 60
	\end{align}
	All elementary cardinalities are nonnegative, so the data are logically consistent. The number of users with none of the three habits is:
	\begin{align}
		x_0 = \card{S} - \sum_{k=1}^7 x_k = 800 - (140+130+60+120+60+90+60) = 800 - 660 = 140.
	\end{align}
\end{solproblema}
